\section{Inference Methods}
\label{sec:methods}
Appendix~\ref{app:implementation} gives decoding settings, molecular scoring details, and failure handling for the following adaptations.

\paragraph{Beam search.}
The deterministic baseline \citep{beam2014} uses width 20 and returns the first sequence in the model's native ordering. All 20 returned sequences contribute to workload, while only rank zero determines success. No property reranking or fallback to a lower-ranked valid molecule is added.

\paragraph{Repeated Sampling with verification.}
The recorded configuration generates 40 proposals and examines their fixed first 20 in original order, implementing sampling with task verification \citep{repeatedsampling2024}. The molecular verifier is exactly the benchmark's chemical eligibility and strict all-target directional-improvement predicate. The first passing proposal is returned; if none passes, proposal zero is returned unchanged. Duplicates and invalid strings retain their positions. No scalar property score or additional tie-breaker is introduced. Final-output SR therefore equals first-20 coverage by construction, rather than demonstrating a learned selector. All 40 generated proposals remain in workload accounting.

\paragraph{Self-Refine.}
The same model generates an answer, critiques it, and revises it \citep{selfrefine2023}. We fix one refinement and return that answer. The recorded prefix contains four responses: initial answer, initial feedback, refinement, and terminal feedback. The last feedback does not change the output. Feedback uses no external candidate properties, and initial and refined answers are not retrospectively compared. This evaluates one feedback--revision iteration, without establishing the benefit of longer trajectories.

\paragraph{Tree of Thoughts.}
ToT retains language-model evaluation that causally determines branches \citep{tot2023}. Five initial molecular states receive five voting responses from the same checkpoint, retaining three parents. Each parent produces five new children; five further voting responses select one terminal child. This gives 20 molecular proposals and ten judging responses. Voting prompts exclude externally computed candidate properties, validity labels, and success labels. Malformed votes contribute zero, and ties follow original frontier order. Molecular prompts, depth, and width are adaptation choices; final selection remains within the terminal frontier without property-based replacement.

\paragraph{Graph of Thoughts.}
G12 instantiates GoT's graph operations \citep{got2024}: 12 initial nodes are scored, six parents retained, and six refinements, six two-parent merges, and four four-parent aggregations generated, totaling 28 nodes. Multi-parent prompts combine molecular design ideas, rather than explicitly editing chemical graphs. Our task-specific scorer combines property improvements, source similarity, synthetic accessibility, and a duplicate penalty. Its features, weights, and molecular feedback are benchmark adaptations. The scorer selects parents and ranks the completed graph; the native full-graph rank-zero node is returned without another reranker.

\paragraph{s1 budget forcing.}
We adapt s1's inference-time budget forcing \citep{s1_2025}, without s1K training or checkpoint replacement. Greedy reasoning permits at most one literal \texttt{Wait} continuation when the initial response ends before its allowance is exhausted. A separate final-answer cue requests the molecule; only that newly generated response supplies the output. This tests an inference component, not the complete trained s1 system.

\paragraph{Exact Voting.}
Four hundred proposals each vote for an answer class \citep{voting2024}. Our molecular equivalence relation groups valid outputs by canonical isomeric SMILES; invalid outputs remain in classes defined by whitespace-normalized response text. The largest class wins, with ties resolved by first appearance and its earliest representative returned. Canonicalization introduces no property-quality preference. Invalid or empty classes can win, without replacement by another candidate.

\paragraph{Four-shot ICL.}
This contextual baseline \citep{brown2020} draws four demonstrations deterministically from the permitted training pool, retaining their task labels. Declared held-out property combinations are excluded. Demonstration source and target molecules are screened against all evaluation sources by raw and canonical identity. No similarity retrieval is added. Four-beam decoding returns its native first beam. The baseline changes prompt context rather than reproducing a complete reasoning-time scaling system.

\paragraph{Reflexion.}
The adaptation retains actor generation, evaluator feedback, failure-conditioned reflection, and within-input memory \citep{reflexion2023}. Each trial generates a thought and molecular action. Property tools inform the evaluator, but model feedback exposes only binary success and invalid-output status, not numerical improvements. Failed nonterminal trials trigger reflection for the next trial. Execution stops on success or after three trials and returns that terminal action, without retrospective best-trial selection. Memory is not shared across inputs.
